\subsection{Prime descent and finite orbit obstructions}
\label{sec:prime-descent}
The following families use the same mechanism. Solvability supplies a
last cyclic-prime extension or equivariantization
\cite[Theorem~1.6 and Proposition~4.4]{ENOSolvable}. When the adjoint
subcategory is the whole category, an extension is impossible. The
remaining equivariantization is constrained by the complete action of
its character subgroup on the source simples. In an equivariantization,
a simple above an orbit of $Y$ has dimension
\[
 [G:G_Y]\dim(\pi)\FPdim(Y),
\]
where $\pi$ is an irreducible projective representation of $G_Y$
\cite[Section~2]{BurciuNatale}. The source-bound checks below include
all dimension products, adjoint closures, and character-subgroup
permutations. Thus descent is a categorical argument applied to
checked hypotheses, rather than an identification by dimension type.

\begin{proposition}\label{prop:primefamilies}
None of the twelve rings
\[
\begin{gathered}
R7\_2\_50,\ R7\_2\_70,\ R7\_2\_113,\ R7\_2\_213,
\ R7\_2\_273,\ R7\_2\_280,\\
R7\_3\_386,\ R7\_3\_398,\ R7\_3\_478,\ R7\_3\_966,
\ R7\_4\_358,\ R7\_4\_963
\end{gathered}
\]
is the Grothendieck ring of a complex fusion category.
\end{proposition}
\begin{proof}
We organize the argument by the descent obstruction.

\emph{Dimensions $40$ and $136$.}
For $R7\_2\_70$, $R7\_2\_213$, $R7\_4\_358$ the type is
$(1^4,2,d,d)$, with $(p,d)=(5,4)$ or $(17,8)$ and global dimension
$8p$. The universal grading is trivial, and every invertible fixes each
noninvertible. Prime descent therefore gives
$\mathcal C=\mathcal D^{C_2}$ with $\mathcal D$ of type
$(1^4,(d/2)^4)$ and dimension $4p$.
If $\mathcal D$ admitted a $C_2$ grading, the outer involution would
preserve a nonzero character of its $2$-primary grading quotient;
this would descend to a nontrivial grading of $\mathcal C$. A $C_p$
grading is also impossible: at $p=17$ its number of components exceeds
rank eight; at $p=5$ its neutral category is pointed of dimension four
and its other four components have rank one. The latter force trivial
pointed associator, while neither
$\operatorname{BrPic}(\operatorname{Vec}_{C_4})$ nor
$\operatorname{BrPic}(\operatorname{Vec}_{C_2^2})$ has an element of
order five. Hence $\mathcal D=\mathcal E^{C_2}$.

Integrality of the descended dimensions and the strong Frobenius
property \cite{ENOSolvable} leave
$\mathcal E$ of type $(1^2,2^{(p-1)/2})$ and dimension $2p$.
For clarity, at dimension ten the alternatives are pointed,
$(1^6,2)$, and $(1^2,2^2)$: the middle case violates subcategory
integrality, while the pointed case gives the forbidden invariant
$C_2$ grading. At dimension thirty-four the dimension partitions are
$(1^2,2^{8-4b},4^b)$, $b=0,1,2$; divisibility excludes $b>0$.
The same grading test gives $\mathcal E_{\rm ad}=\mathcal E$, so
$\mathcal E=(\operatorname{Vec}_{C_p}^{\omega})^{C_2}$ with inversion
on $C_p$.

The kernel of
$\operatorname{Inv}Z(\mathcal E)\to\operatorname{Inv}\mathcal E$ is
$\operatorname{Aut}_{\otimes}(\operatorname{Id}_{\mathcal E})
=\operatorname{Hom}(U(\mathcal E),\C^\times)=1$.
Consequently its central $\operatorname{Rep}(C_2)$ is unique and is
preserved by the successive outer actions. Combining the
 equivariantizations \cite{DGNO,ENOSolvable} gives
$\mathcal C=(\operatorname{Vec}_{C_p}^{\omega})^H$, where $|H|=8$ and
$H_{\rm ab}=C_2^2$. Its action on $C_p$ has cyclic image of order at
most two. At $p=17$, orbit size times projective stabilizer degree is
at most four, contradicting the dimension-eight simples.
At $p=5$, choose a strictly $H$-invariant $\mu_p$-valued associator by
averaging. The vanishings $H^2(C_p,\C^\times)=0$ and
$H^2(H,\widehat{C_p})=0$ remove the action tensorators and their
coherence multipliers. Thus the stabilizer representations are ordinary;
their degrees are at most two, and nontrivial orbits have abelian
stabilizer of order four. All dimensions are at most two, contradicting
$d=4$.

\emph{Dimension $112$.}
For $R7\_3\_386$, $R7\_3\_966$, type $(1^2,2^3,7^2)$ and trivial
universal grading force a $C_2$ descent of type $(1^7,7)$ and dimension
$56$. Its adjoint category is again the whole category. Its last prime
step is therefore a $C_7$ equivariantization of
$\operatorname{Vec}_{C_2^3}^{\omega}$, with a Singer seven-cycle on
the nonzero elements. The unique central $C_7$ survives the outer
involution, so the combined group has order fourteen. Its representation
subcategory has type $(1^2,2^3)$ and hence is $\operatorname{Rep}(D_{14})$.
The resulting action would conjugate the Singer cycle to its inverse.
Its normalizer in $\operatorname{GL}_3(\mathbb F_2)$ is $7{:}3$, which
contains no such involution. This is impossible.

\emph{Dimension $96$.}
For $R7\_2\_113$, $R7\_2\_273$, the character involution fixes the
dimension-two and both dimension-six objects and exchanges the two
dimension-three objects. These are exactly the premises of the
dimension-$96$ proof above; its $(1,3,6,6)$ orbit contradiction applies.
For $R7\_2\_50$, $R7\_2\_280$, the involution instead exchanges both
the dimension-three pair and the dimension-six pair. Descent has type
$(1^3,3,6)$ and dimension $48$, with trivial universal grading. Its
next descent is a $C_3$ equivariantization of a dimension-$16$ category
$\mathcal E$ of type $(1^4,2^3)$, with the three nontrivial invertibles
and the three dimension-two objects each cyclically permuted.
The invertibles form $V=C_2^2$ and fix the dimension-two objects.
Hence $\mathcal E_{\rm ad}=\operatorname{Vec}_V^{\omega}$, with a
faithful $V$ grading whose three nonneutral components have rank one.
A rank-one pointed module implies $[\omega]=0$ by its module pentagon.
There are only two rank-one left module categories over
$\operatorname{Vec}_V$, since $H^2(V,\C^\times)=C_2$.
Distinct invertible grading components here must have inequivalent left
module categories: an equivalence, tensored with the inverse of one
component, would identify a nonneutral rank-one component with the
rank-four regular module. Three such components cannot fit into two
classes.

\emph{Dimension $54$.}
For $R7\_3\_398$, $R7\_3\_478$, type $(1^2,2^4,6)$ forces descent
of type $(1^9,3^2)$ and dimension $27$. Its pointed adjoint has group
$A$ of order nine; the other two $C_3$ grading components have rank one,
so its associator is trivial. The outer involution inverts both $A$ and
the grading $C_3$. Its restricted tensorator has class
$\eta\in H^2(A,\C^\times)$. Inversion acts trivially on this
odd-order group, while order-two coherence gives $2\eta=0$; hence
$\eta=0$. Its Brauer--Picard class restricts to $-I$ on
$A\oplus\widehat A$, which is central in the orthogonal group
\cite[Corollary~1.2]{ENOExt}. It therefore cannot conjugate a nontrivial
order-three component class to its inverse, as the grading action
requires.

\emph{Dimension $310$.}
For $R7\_4\_963$, type $(1^2,2^2,10^3)$ has dimension
$310=2\cdot5\cdot31$. The $pqr$ theorem
\cite[Theorem~9.2]{ENOSolvable} makes any categorification
$\mathcal C(G,\omega,H,\psi)$ with $|G|=310$.
The complete closed-subring scan gives dimensions $1,2,10,310$, so
$|H|$ belongs to this list. For $|H|=1,2$ the double-coset formula
\cite[Theorem~5.1]{GelakiNaidu} cannot supply dimension ten.
For $|H|=10$, the representation subcategory forces $H=D_{10}$.
The three outside dimension-ten simples force trivial double-coset
intersections and outside cosets of size one hundred. Thus $G$ would
be a Frobenius group with kernel $C_{31}$ and complement $D_{10}$;
a faithful complement action in the cyclic group
$\operatorname{Aut}(C_{31})$ is impossible.
For $H=G$, the category is $\operatorname{Rep}(G)$. Its normal Sylow
$31$ subgroup has a split complement $D_{10}$, whose cyclic action on
$C_{31}$ has image at most two. The orbit/stabilizer formula then
bounds all representation degrees by two, again contradicting ten.
\end{proof}

\subsection{Four additional exhaustive signed exclusions}
The complete oriented search excludes
$R7\_3\_1036$, $R7\_3\_1044$, $R7\_4\_2510$, $R7\_4\_2706$.
Their complete lists have respectively $8,24,4,48$ signed associative
tables. Every possible dimension character fails either nonvanishing,
duality, or formal-codegree norm divisibility. The supplement regenerates
the lists and rechecks each branch, including under optimized Python.
These are general exclusions: sphericalization doubles both global
dimension and the relevant formal codegrees, leaving the necessary
quotient unchanged. No pivotal structure on the original category is
assumed.

\subsection{A relative-trace obstruction without signed enumeration}
\begin{proposition}[Relative-trace criterion]
\label{prop:relative-trace}
Let $R$ be a commutative fusion ring with self-dual basis, and put
$B=\sum_{i,j,k}(N_{ij}^k)^2$. Suppose $L\in\Z_{>0}$ and $f$ are
formal codegrees of $R$, and set $F=\mathbb Q(f)$.
If every nonzero totally positive element of
$L\mathcal O_F\cap f\mathcal O_F$ has field trace greater than $B$,
then $R$ has no categorification over $\C$.
\end{proposition}
\begin{proof}
Sphericalize a hypothetical category and take its quotient by
$\varepsilon=-1$. Both lifts of a self-dual simple remain self-dual
\cite[Section~5.1(a)]{OstrikRank3}; hence the signed multiplication
matrices are symmetric and commute. Writing them as $M_i$, their
Casimir satisfies
\[
 \operatorname{Tr}\Bigl(\sum_iM_i^2\Bigr)
 =\sum_{i,j,k}(T_{ij}^k)^2\le B.
\]
Let $d_i$ be the lift dimensions, $K=\mathbb Q(d_i)$, and
$\delta=\sum_i d_i^2$. Every embedding of $K$ gives a distinct
dimension character and contributes its conjugate of $\delta$ as
a Casimir eigenvalue. This counts $[K:\mathbb Q]$ characters even
when some eigenvalues coincide. All Casimir eigenvalues are positive,
so $\operatorname{Tr}_{K/\mathbb Q}(\delta)\le B$.

The sphericalization has dimension $2\delta$ and formal codegrees
$2L,2f$. The induced-unit theorem gives algebraic integers
$\delta/L,\delta/f$ in its dimension field $K$
\cite[Theorem~2.13 and Corollaries~2.14--2.15]{OstrikRank3}.
Thus $F\subseteq K$. The element
$t=\operatorname{Tr}_{K/F}(\delta)$ is totally positive and lies in
$L\mathcal O_F\cap f\mathcal O_F$. Its trace is at most $B$, a
contradiction. No equality $K=\mathbb Q(\delta)$ is assumed.
\end{proof}

\begin{proposition}
The rings $R7\_4\_1417$ and $R7\_4\_1420$ admit no complex
fusion-category categorification.
\end{proposition}
\begin{proof}
Both rings have formal codegrees $20$ and
$f=510+70\sqrt{51}$, whose minimal polynomial is
$z^2-1020z+10200$. Here
$\mathcal O_F=\Z[\sqrt{51}]$. Write
$t=20(x+y\sqrt{51})$ in the intersection ideal. Then
\[
 \frac{t}{f}=(x-7y)+\left(y-\frac{7x}{51}\right)\sqrt{51}.
\]
Integrality forces $51\mid x$, while total positivity forces $x>0$.
Consequently $\operatorname{Tr}_{F/\mathbb Q}(t)=40x\ge2040$.
The exact Casimir trace bounds are $1086$ and $1070$, respectively,
so Proposition~\ref{prop:relative-trace} applies. The source-bound
checker recomputes both characteristic polynomials, trace bounds,
and the displayed quotient identity using exact arithmetic.
\end{proof}


\subsection{Pivotal exclusions within the general frontier}
The rings $R7\_4\_1212$ and $R7\_4\_1213$
have no pivotal categorification. The complete ordinary character
screen excludes every candidate dimension character either by a zero
simple dimension or by nonintegrality of $\delta/f$, where
$\delta=\sum_i d_i d_{i^*}$ and $f$ is an ordinary formal codegree
\cite[Corollary~2.14]{OstrikRank3}. The exact character branches and
norm witnesses are separately certified. This assertion excludes
spherical, pseudo-unitary and unitary categories as well, but leaves
general complex categorification unresolved. These two rings remain
among the $39$ undecided entries in Table~\ref{tab:status}; no pivotal
obstruction is counted as a general exclusion.
